% © 2025 Ing. David Jaroš — CC BY-NC-ND 4.0
%
% This work is licensed under a Creative Commons Attribution-NonCommercial-NoDerivatives
% 4.0 International License (CC BY-NC-ND 4.0).
%
% Biquaternion Algebra — Canonical Definition
% Status: Canonical
% Version: 1.0
% Date: 2026-03-01

\section{Biquaternion Algebra $\mathcal{B} = \mathbb{C}\otimes\mathbb{H}$}
\label{sec:algebra:biquaternion_algebra}

\subsection{Definition}

The \textbf{biquaternion algebra} $\mathcal{B}$ is the tensor product of the field
of complex numbers with the quaternion algebra:

\begin{equation}
  \label{eq:biquat_def}
  \mathcal{B} = \mathbb{C} \otimes_{\mathbb{R}} \mathbb{H}.
\end{equation}

A general biquaternion has the form
\begin{equation}
  \label{eq:biquat_general}
  q = z_0 + z_1\,\mathbf{I} + z_2\,\mathbf{J} + z_3\,\mathbf{K},
  \qquad z_\mu \in \mathbb{C},
\end{equation}
where $\{1, \mathbf{I}, \mathbf{J}, \mathbf{K}\}$ is the standard quaternion basis
satisfying
\begin{equation}
  \label{eq:quat_rules}
  \mathbf{I}^2 = \mathbf{J}^2 = \mathbf{K}^2 = \mathbf{I}\mathbf{J}\mathbf{K} = -1.
\end{equation}

\subsection{Matrix Isomorphism}

There is a canonical $\mathbb{C}$-algebra isomorphism
\begin{equation}
  \label{eq:biquat_iso}
  \mathcal{B} \;\cong\; \mathrm{Mat}(2, \mathbb{C}),
\end{equation}
given explicitly by
\begin{equation}
  1 \mapsto I_2, \quad
  \mathbf{I} \mapsto i\sigma_1, \quad
  \mathbf{J} \mapsto i\sigma_2, \quad
  \mathbf{K} \mapsto i\sigma_3,
\end{equation}
where $\sigma_1, \sigma_2, \sigma_3$ are the Pauli matrices.

\subsection{Involutions and $\mathbb{Z}_2 \times \mathbb{Z}_2 \times \mathbb{Z}_2$}

The imaginary subspace $\mathrm{Im}(\mathbb{H}) = \mathrm{span}_{\mathbb{R}}\{\mathbf{I},
\mathbf{J}, \mathbf{K}\}$ admits three commuting involutions:
\begin{align}
  P_{\mathbf{I}}(q) &= -\mathbf{I}\,q\,\mathbf{I}^{-1}, \\
  P_{\mathbf{J}}(q) &= -\mathbf{J}\,q\,\mathbf{J}^{-1}, \\
  P_{\mathbf{K}}(q) &= -\mathbf{K}\,q\,\mathbf{K}^{-1}.
\end{align}
These satisfy $P_{\mathbf{I}}^2 = P_{\mathbf{J}}^2 = P_{\mathbf{K}}^2 = \mathrm{id}$
and generate $\mathbb{Z}_2 \times \mathbb{Z}_2 \times \mathbb{Z}_2$.
Their Lie algebra equals $\mathfrak{su}(3)$; see
\texttt{canonical/algebra/involutions\_Z2xZ2xZ2.tex} and
\texttt{canonical/algebra/algebra\_summary\_table.tex}
for the detailed derivation.

\subsection{Automorphism Group}

\begin{proposition}[Automorphism group of $\mathcal{B}$]
  \label{prop:aut_B}
  $\mathrm{Aut}_{\mathbb{R}}(\mathcal{B}) \cong
  \bigl[\mathrm{GL}(2,\mathbb{C}) \times \mathrm{GL}(2,\mathbb{C})\bigr]/\mathbb{Z}_2.$
\end{proposition}

The left factor of $\mathrm{GL}(2,\mathbb{C})$ gives rise to $\mathrm{SU}(2)_L \times
\mathrm{U}(1)_Y$ (electroweak gauge group) after restriction to unitary elements.
The right factor gives $\mathrm{U}(1)_{\mathrm{EM}}$.

\subsection{Key Algebraic Properties}

\begin{itemize}
  \item \textbf{Associative}: $\mathcal{B}$ is associative (inherited from
    $\mathbb{H}$ and $\mathbb{C}$, or via $\mathrm{Mat}(2,\mathbb{C})$).
  \item \textbf{Non-commutative}: $\mathbf{I}\mathbf{J} \neq \mathbf{J}\mathbf{I}$
    in general.
  \item \textbf{Dimension}: $\dim_{\mathbb{R}} \mathcal{B} = 8$;
    $\dim_{\mathbb{C}} \mathcal{B} = 4$.
  \item \textbf{Centre}: $Z(\mathcal{B}) = \mathbb{C} \cdot 1$ (complex scalars).
\end{itemize}

\subsection{Role in UBT}

The fundamental field of UBT is a $\mathcal{B}$-valued section
$\Theta(q, \tau)$ over biquaternion coordinate $q$ and complex time
$\tau = t + i\psi$.  All physical content — spacetime geometry, gauge symmetries,
and particle content — emerges from $\Theta$ and the algebraic structure of
$\mathcal{B}$.

The three-dimensionality $\dim_{\mathbb{R}}(\mathrm{Im}\,\mathbb{H}) = 3$ directly
forces $\mathrm{SU}(3)_c$ as the color gauge group with zero free parameters.

% See also:
%   canonical/algebra/algebra_summary_table.tex   — quick-reference table
%   canonical/algebra/involutions_Z2xZ2xZ2.tex    — Z2xZ2xZ2 involution structure
%   canonical/fields/biquaternion_algebra.tex      — detailed field-theory treatment
